\ifx\wholebook\relax \else

\documentclass[b5paper]{article}
\usepackage[nomarginpar
  %, margin=.5in
]{geometry}

\addtolength{\oddsidemargin}{-0.05in}
\addtolength{\evensidemargin}{-0.05in}
\addtolength{\textwidth}{0.1in}

\usepackage[en]{../prelude}

\setcounter{page}{1}

\begin{document}

\title{序}

\author{刘新宇
\thanks{{\bfseries 刘新宇} \newline
  Email: liuxinyu99@hotmail.com \newline}
  }

\maketitle
\fi

\markboth{Forward}{The Journey of Numbers}

\chapter*{Forward}
\phantomsection  % so hyperref creates bookmarks
\addcontentsline{toc}{chapter}{Forward}

I had a chance to listen to my kid's class during the Covid-19 pandemic in 2022. The entire city was locked down, and we had to stay at home, working remotely and learning remotely without knowing when it would be an end. It was a hard time for all, in particular, for teachers. They tried their best to continue the courses over the winter semester although all were infected and in bad condition. Without face to face interaction, without effective feedback, it's hard; it's all about insisting.

As a father, what can I help? People discovered more and more knowledge age by age. What we teach to our children is highly summarized. For example, the mathematics taught is in a state of done. There were great minds and interesting stories behind each theorem: like Pierre de Fermat's `secrete study', the knightly challenge between Cardano and Tartaglia for the cubic equations, the young man who discovered geometric meaning of complex numbers but forgot to leave his name... Those beautiful result arose from failures, errors, tears, and joys, but we don't have enough time in the class to tell them to our children. This is exactly what our parents may help: when the pandemic is over, visit the museum with our kids to understand the history, the people, and the `why' for those beautiful achievements.

Mathematics is about and arose from pursuing truth and beauty. Everything has to be justified rigorously. There is no authority; great mathematicians, like Fermat and Euler made mistakes. It is through the continuously identifying and correcting defects that mathematics progress. There were many times in history, that people completely rebuilt mathematics on top of a genuine new foundation. Many mathematicians were true scholars. When the young Lagrange, at the age of 19 who was unknown to the public, wrote Euler with overlapping discoveries in the calculus of variations, Euler decided to conceal his result until Lagrange publishing first, ``in no way do I want to take away from you any of the glory that you deserve.'' When Emmy Noether was facing obstacle of being a faculty member only because she was a woman, Hilbert fought for her and say, ``this is a University but not a bathhouse.''

她追求美，简洁的美，对称的美，清晰的美，抽象的美。克罗内克在博士答辩时坚持认为数学是艺术；埃尔德什一生云游追寻“数学天书”中最美的证明；欧拉公式被称作“最美公式”；黄金分割数指导着艺术家们对于美的创作；迷人的分形图案出现在大自然中。

在课堂教学的主阵地之外摇旗呐喊，带孩子们走进数学博物馆研学，体验数学的历史和真善美，就是我写作本书的目的。这本书不是课本，不和课堂教学冲突；它不是教辅，不拆解例题，不讲解题技巧；它不是刷题手册，尽管每章带有动手实践的习题和完整的答案；它不是通俗科普，不讲传说八卦。这是一本认真的数学科普书，以数的发展为线索，带着青少年读者沿着历史脉络“参观”数学博物馆。回答为什么负负为正？为什么零不能作除数？为什么有的分数化成小数会循环？为什么圆周率在小数点后无穷无尽？……这本书同时回答“怎样做？”的问题，古人怎样计算圆周率？怎样开方？怎样用尺规作出正五边形，数学家怎样找出三次方程的求根公式……我们刨根问底，分数线是谁发明的？小数点是谁发明的？百分号是谁引入的？央行降息的基点是啥意思？为什么自然对数的底叫做e？我们破除谣言，毕达哥拉斯不可能听到铁匠铺敲打出不同音高的和声；西帕索斯因发现无理数被谋杀只是个传说；书本纸张并非是黄金比的长方形……

这本书的前一半只涉及义务教育阶段的数学，后一半涉及高中阶段的数学。其中圆周率在祖冲之以后的计算方法涉及了高中的微积分初步知识。此外从第4章开始，加入了初等数论的内容。毕竟这是关于数的科普书，数论是绕不过去的。

本书成书过程中得到了常成的大力帮助，仔细审阅了我的初稿。机械工业出版社的编辑们给了我不断的支持和鼓励。家里的大人孩子和很多朋友也成了我的义务审稿人，在这里一并感谢。

\vspace{15mm}

刘新宇

二零二五年深秋于北京和平里

%% 本书的电子版可以在\url{https://github.com/liuxinyu95/amathtour}下载。如果希望获得纸质版，请联系作者：liuxinyu99@hotmail.com

\ifx\wholebook\relax \else

\expandafter\enddocument
%\end{document}

\fi
